\chapter*{Introduction}
It was the 4 July 2012 when Joseph Incandela, CMS spokesperson, affirmed: "\textit{They combine to give us a combined significance of five standard deviations}". With these exact words, he announced, confirmed by ATLAS spokesperson Fabiola Gianotti a few minutes later, the first observation of the Higgs boson. Twelve years have passed since then and it remains the last big breakthrough in High Energy Physics, completing the Standard Model (SM) framework. What happened after that? Has Particle Physics come to an end and everything that needed to be discovered has been found? \\
In reality this field of research is as active as it was a few years back, currently pursuing precision's measurements and search for deviations from the Standard Model, such as Dark Matter. In particular, this second kind of research has been conducted using traditional data analysis techniques targeting specific theoretical models but, unfortunately, nothing has been observed yet at collider experiments and only upper limits on production cross-sections were set. There are many reasons for this unexpected outcome: wrong Physical model for experiments, wrong triggering while selecting data, or deviations too small to be observed through normal techniques.\\
A possible solution could be Anomaly Detection, that is the search for data deviating from what is already known. In our field of interest, Particle Physics, we apply this method to look for anomalies (signal data), deviating from well known reacting and decaying particles (background data). The strength of this method is that it is model agnostic, meaning that it does not rely on a specific theoretical model to analyse data. This allows us to explore a more vast range of events and gives us a first hint of the region where to look more in deep for New Physics with precise data analysis techniques. In this process, Machine Learning (ML) and Deep Learning architectures (DL) come to help and speed up the process.\\




An example case has been investigated in this thesis in order to understand the potential of this tool. Using Monte Carlo generated data as sample, anomalies in particles decay have been studied, implementing a DL architecture, in particular a Transformer. This thesis has been an opportunity to dive into this topic and become confident with tools used in this field of research.\\
The thesis is structured as follows. Chapter \ref{chap:LHCAtlas} introduces the LHC and ATLAS, experiments to investigate particle collisions. Chapter \ref{chap:AD} presents Anomaly Detection approach. In chapter \ref{chap:Dataset} the dataset analysed is presented. In chapter \ref{chap:Model} is possible to find a detailed description of the architecture implemented. In conclusion, chapter \ref{chap:Results} shows results of this project.\\\\

\noindent \textbf{First index draft:}
\begin{itemize}
    \item Introduction
    \item SM \& BSM
    \begin{itemize}
        \item Particles
        \item QED
        \item QCD
        \item EW
        \item Higgs and particles collisions
        \item BSM
    \end{itemize}
    \item ATLAS
    \begin{itemize}
        \item LHC accelerator system
        \item ATLAS detector
    \end{itemize}
    \item Anomaly Detection
    \begin{itemize}
        \item What is it?
        \item How is it implemented using DNNs?
        \item Why AD in our case? How are we implementing it?
    \end{itemize}
    \item How is a Particle Physics Analysis conducted?
    \begin{itemize}
        \item MC sample
        \item distributions
        \item ...
        \item CLs method
    \end{itemize}
    \item ADJJ analysis: plots and results
\end{itemize}



